Pour la première partie de nos travaux, à savoir le rejeu d'une spécification
paramétrique, nous avons plusieurs objectifs à atteindre. %

Le premier est l'enregistrement d'un processus de modélisation. %
Pour que cela soit possible, il faut construire un nom persistant (NP) pour
chaque entité topologique du modèle et enregistrer dans la spécification
paramétrique les entités qui doivent être passées en paramètres topologiques des
opérations. %
Le deuxième objectif est le suivi des changements topologiques. %
Les opérations étant formalisées par des règles avec une syntaxe, celles-ci
décrivent précisément et localement les changements topologiques à l'issue de
leur application. %
Le suivi de ces changements topologiques est nécessaire pour la construction
d'un journal de bord, une structure de donnée dont le but est d'enregistrer tous
les changements topologiques, pour toutes les orbites, appliquées par toutes les
opérations. %
Le troisième objectif est le rejeu. %
Cette étape nécessite, dans un premier temps, de construire dans un graphe
orienté acyclique (ou DAG) de jeu, pour chaque paramètre topologique, son
historique de changement topologique. %
Dans un second temps, de construire un DAG de rejeu pour effectuer l'appariement
des entités topologiques et les retrouver dans le nouveau modèle. %
Le quatrième objectif, qui conclut cette partie, est le support des scripts de
règle. %

Pour la deuxième partie, l'objectif et d'animer l'évolution d'un modèle sur une
échelle de temps. %
Cet objectif peut comprendre, d'une part, l'intégration de paramètres temps dans
les opérations pour simuler, par exemple, l'érosion sur une surface rocheuse. %
D'une autre part, cela peut comprendre l'insertion d'opération à des temps
donnés pour représenter de manière spontanée des événements historiques et leur
impactes sur des bâtiments. %

Enfin, pour la troisième partie, l'objectif est d'intégrer des données
historiques dans les modèles. %
Pour cela, l'outil final devra pouvoir stocker divers types de données et les
présenter lorsque l'utilisateur sélectionnera tout ou une partie d'un modèle. %
Par exemple, en sélectionnant un pan de mur, l'utilisateur pourra avoir accès à
des informations telles que la date de construction du bâtiment, voire du mur
s'il s'agit d'une extension, du matériau utilisé, de sa rugosité ou encore de
son niveau d'érosion. %


%%% Local Variables:
%%% mode: latex
%%% TeX-master: "../main"
%%% End:
